% Generated draft: requires mathematical and source review.
\documentclass[11pt]{article}
\usepackage[margin=25mm]{geometry}
\usepackage{fontspec}
\setmainfont{Latin Modern Roman}
\usepackage{amsmath,amssymb,mathtools,unicode-math}
\setmathfont{Latin Modern Math}
\usepackage{enumitem,needspace,xurl,hyperref}
\setlength{\parindent}{0pt}
\setlength{\parskip}{5pt}
\setlength{\emergencystretch}{3em}
\newcommand{\N}{\mathbb{N}}
\newcommand{\Z}{\mathbb{Z}}
\newcommand{\Q}{\mathbb{Q}}
\newcommand{\R}{\mathbb{R}}
\newcommand{\C}{\mathbb{C}}
\newcommand{\F}{\mathbb{F}}
\newcommand{\A}{\mathbb{A}}
\newcommand{\PP}{\mathbb{P}}
\newcommand{\E}{\mathbb{E}}
\newcommand{\eps}{\varepsilon}
\newcommand{\dd}{\,\mathrm{d}}
\begin{document}
\section*{UnsolvedMath ID: 701416}
\textbf{Source problem code:} catalogue entry
\subsection{Definitions and mathematical statement}
A group \(G\) is locally free if every finitely generated subgroup of \(G\) is a free group. No finite generation assumption is imposed on \(G\) itself. A maximal proper subgroup is a subgroup \(M<G\) for which there is no subgroup \(H\) with \(M<H<G\). Define the Frattini subgroup by
\[
\Phi(G)=\bigcap_{\substack{M<G\\ M\text{ maximal proper}}}M,
\]
with the intersection interpreted as \(G\) when there are no maximal proper subgroups.

\textbf{Question (J. Wiegold).} Does there exist a group \(G\) such that every finitely generated subgroup of \(G\) is free, there exist \(a,b\in G\) with \(ab\ne ba\), and
\[
\Phi(G)\ne\{1_G\}?
\]
The supplied record labels this question open; that label is not an independently verified mathematical or literature assessment.
\subsection{Short English statement}
Can a nonabelian group whose every finitely generated subgroup is free have a nontrivial Frattini subgroup?
\subsection{Sources}
\begin{itemize}
\item Local UnsolvedMath registry: permanent ID \(701416\), exact source problem code \texttt{catalogue entry}; supplied upstream revision \texttt{d371ee917ac9647fbb06cd6fc529deec6f3dda82}.
\item J. Wiegold, Problem 17.115, \emph{The Kourovka Notebook (Unsolved Problems in Group Theory)}, 17th issue (2010), as attributed in the supplied record.
\item The supplied background also identifies the 21st issue edition (2026), arXiv:1401.0300v48. The supplied collection description names E. I. Khukhro and V. D. Mazurov as editors. These bibliographic details have not been independently verified.
\end{itemize}
\end{document}
